% 主 LaTeX 檔案：main.tex
% 這是教材的總檔案，使用 \include{} 來包含各章節
% 您可以新增其他章節檔案，如 chapter2.tex 等

\documentclass[12pt,a4paper]{book}
\usepackage[utf8]{inputenc}
\usepackage{geometry}
\geometry{left=2.5cm,right=2.5cm,top=2.5cm,bottom=2.5cm}
\usepackage{fancyhdr}
\usepackage{graphicx}
\usepackage{amsmath}
\usepackage{hyperref}
\usepackage{verbatim}

% 這是一份latex 文件的 preamble
% 請AI閱讀後，將附上的PDF整理，以latex code整理，儘量用longtable
% 請注意我的名字: 謝慕揚, MD, PhD, FESC
% 表格請注意換行，不該在cell內使用 \\
% 表格請不要有垂直分隔線 |



% Including essential packages for Chinese typesetting
\usepackage{xeCJK}
\usepackage{fontspec}
% Configuring Chinese font with Noto Serif CJK TC, fallback to SimSun
\IfFontExistsTF{Noto Serif CJK TC}{
  \setCJKmainfont{Noto Serif CJK TC}
  \setCJKsansfont{Noto Serif CJK TC}
  \setCJKmonofont{Noto Serif CJK TC}
}{\setCJKmainfont{SimSun}
  \setCJKsansfont{SimSun}
  \setCJKmonofont{SimSun}
}

% Including other necessary packages
\usepackage{geometry}
\usepackage{titlesec}
\usepackage{enumitem}
\usepackage{xcolor}
\usepackage{colortbl}
\usepackage{tcolorbox}
\usepackage{booktabs}
\usepackage{longtable}
\usepackage{array}
\usepackage{graphicx}
\usepackage{tikz}
\usepackage{fancyhdr}
\usepackage{multicol}
\usepackage{datetime2}
\usepackage{hyperref}
\usepackage{mdframed}
\usepackage{amsmath}

\hypersetup{
    colorlinks=true,
    linkcolor=emergency_blue,
    citecolor=emergency_blue,
    urlcolor=emergency_blue,
    pdfborder={0 0 0}
}

% Defining page geometry
\geometry{
  top=2.5cm,
  bottom=2.5cm,
  left=2cm,
  right=2cm,
  headheight=25pt
}

% Adjusting topmargin to compensate for increased headheight
\addtolength{\topmargin}{-10pt}

% Defining colors
\definecolor{emergency_red}{RGB}{186,24,27}
\definecolor{emergency_blue}{RGB}{0,114,188}
\definecolor{emergency_gray}{RGB}{120,120,120}
\definecolor{highlight_yellow}{RGB}{255,250,205}

% Setting title formats
\titleformat{\section}[block]{\Large\bfseries\color{emergency_red}}{\thesection}{10pt}{}
\titleformat{\subsection}[block]{\large\bfseries\color{emergency_blue}}{\thesubsection}{8pt}{}
\titleformat{\subsubsection}[block]{\normalsize\bfseries\color{emergency_gray}}{\thesubsubsection}{6pt}{}

% Defining custom environment for important notes
\newmdenv[
    linecolor=emergency_red,
    backgroundcolor=highlight_yellow,
    linewidth=2pt,
    topline=true,
    bottomline=true,
    leftline=true,
    rightline=true,
    innertopmargin=10pt,
    innerbottommargin=10pt,
    innerleftmargin=10pt,
    innerrightmargin=10pt
]{importantbox}

% Defining note command with tcolorbox
\newcommand{\note}[1]{
    \par
    \noindent
    \begin{tcolorbox}[
        colback=emergency_blue!5!white,
        colframe=emergency_blue,
        title=\textbf{重點提醒},
        fonttitle=\bfseries,
        boxrule=0.5pt,
        arc=3pt,
        left=6pt,
        right=6pt,
        top=6pt,
        bottom=6pt
    ]
    #1
    \end{tcolorbox}
}

% % Setting up header and footer
% \pagestyle{fancy}
% \fancyhf{}
% \fancyhead[L]{\textcolor{emergency_red}{心血管中心教學文件}}
% \fancyhead[R]{\textcolor{emergency_blue}{急性心肌梗塞早期診斷與心電圖教學}}
% \fancyfoot[C]{\textcolor{emergency_gray}{\thepage}}
% \renewcommand{\headrulewidth}{1pt}
% \renewcommand{\headrule}{\hbox to\headwidth{\color{emergency_red}\leaders\hrule height \headrulewidth\hfill}}

% Defining custom date format for ISO 8601
\DTMnewdatestyle{iso}{
  \renewcommand{\DTMdisplaydate}[4]{\number##1-\number##2-\number##3}
  \renewcommand{\DTMDisplaydate}[4]{\number##1-\number##2-\number##3}
}
\DTMsetdatestyle{iso}
\newcommand{\mydate}{\DTMtoday}

\title{實證醫學教材}
\author{黃沛銓 MD\\謝慕揚 MD, PhD, FESC\\黃國良}
\date{\today}

\begin{document}

\maketitle

\tableofcontents

\chapter{實證醫學概述}
% 第一章檔案：chapter1.tex
% 這是第一章的獨立內容，可直接 \include 到主檔案

\section{定義與原則}

實證醫學（Evidence-Based Medicine, EBM）被廣泛定義為「謹慎、明確且明智地使用當前最佳證據，來做出關於個別患者照護的決策」。這一定義源自David Sackett等人於1996年的經典描述，強調實證醫學並非僅依賴傳統經驗或專家意見，而是將最佳可用外部臨床證據與個別臨床專業知識相整合。

實證醫學的核心原則包括以下三個要素：
\begin{itemize}
    \item \textbf{最佳研究證據}：來自系統性研究的臨床相關證據，例如隨機對照試驗（RCT）、系統回顧與meta-analysis。
    \item \textbf{臨床專業知識}：醫師透過臨床經驗與實踐所累積的專業判斷能力，用以有效診斷並考量患者個別情況。
    \item \textbf{患者價值與偏好}：決策必須尊重患者的權利、期望與個人情境，實現共享決策（shared decision-making）。
\end{itemize}

實證醫學的實踐通常遵循五個步驟（Sackett模型）：
\begin{enumerate}
    \item 將臨床問題轉化為可回答的結構化問題（例如使用PICO框架）。
    \item 搜尋最佳可用證據。
    \item 批判性評估證據的有效性、影響力與適用性。
    \item 將證據應用於臨床實踐，並整合臨床專業與患者價值。
    \item 評估實踐成效並持續改進。
\end{enumerate}

這些原則確保醫療決策基於科學證據，而非直覺或陳規，從而提升照護品質與患者安全。

\section{EBM的歷史發展}

實證醫學的哲學起源可追溯至19世紀中葉的巴黎醫學派，以及更早的醫學統計應用。然而，現代實證醫學的正式發展始於20世紀後期。

關鍵歷史里程碑包括：
\begin{itemize}
    \item \textbf{1972年}：英國流行病學家Archie Cochrane出版《Effectiveness and Efficiency》，批評醫療資源分配缺乏隨機對照試驗支持，奠定Cochrane Collaboration的基礎。
    \item \textbf{1980年代}：加拿大McMaster大學的David Sackett領導臨床流行病學部門，發展批判性評估方法，並出版系列文章指導醫師閱讀文獻。
    \item \textbf{1990年}：Gordon Guyatt首次使用「evidence-based medicine」一詞，描述新的醫學教學方法。
    \item \textbf{1992年}：Evidence-Based Medicine Working Group發表五步驟模型。
    \item \textbf{1996年}：Sackett等人於《BMJ》發表經典論文，明確定義EBM並澄清其非「食譜式醫學」。
    \item \textbf{1990年代後期}：成立Cochrane Collaboration與Oxford Centre for Evidence-Based Medicine，推動系統回顧與全球合作。
\end{itemize}

實證醫學從個別決策擴展至指南制定與政策，現已融入醫學教育與臨床實踐。

\section{EBM在醫療決策中的角色}

實證醫學在醫療決策中扮演核心角色，它橋接科學證據與臨床實務，確保決策基於可靠證據而非主觀判斷。這有助於：
\begin{itemize}
    \item 提升決策品質：減少無效或有害介入，促進有效治療。
    \item 資源最適化：避免不必要檢查或治療，降低醫療成本。
    \item 患者中心照護：整合患者價值，促進共享決策，提高滿意度與依從性。
    \item 終身學習：鼓勵醫師持續更新知識，應對醫學快速進展。
\end{itemize}

在臨床情境中，EBM適用於診斷、治療、預後與預防等領域。例如，面對患者問題時，醫師可使用EBM步驟快速搜尋並應用最新證據，取代過時經驗。儘管面臨時間限制或證據不足等挑戰，EBM仍是現代醫學的基石，持續推動醫療品質改進。

% 可在此新增圖表、參考文獻或習題
% \bibliographystyle{plain}
% \bibliography{references} % 若有參考文獻檔案



\chapter{醫療問題的制定}
% 第二章檔案：chapter2.tex
% 這是第二章的獨立內容，可直接 \include 到主檔案的 \chapter{醫療問題的制定} 之後
% 假設主檔案中已有 \chapter{醫療問題的制定}，本檔案從 section 開始

\section{PICO框架（Patient, Intervention, Comparison, Outcome）}

PICO框架是實證醫學中用於制定清晰、可回答臨床問題的核心工具，由Richardson等人於1995年提出，後經Sackett等人在實證醫學教學中廣泛應用。該框架將臨床問題結構化為四個主要組成部分，有助於精準定義問題並指導後續文獻搜尋。

PICO各元件定義如下：
\begin{itemize}
    \item \textbf{P（Patient/Population/Problem）}：描述目標患者或人群，包括疾病、年齡、性別、共病或其他相關特徵。例如：「65歲以上社區居住的老年人」或「患有2型糖尿病的成人」。
    \item \textbf{I（Intervention）}：指主要介入措施，例如治療、診斷測試、暴露因素或預防策略。例如：「有氧運動訓練」或「低劑量阿斯匹靈」。
    \item \textbf{C（Comparison）}：對照組或替代介入，包括標準治療、安慰劑或無介入。例如：「標準藥物治療」或「無運動介入」。
    \item \textbf{O（Outcome）}：預期測量的結果，包括主要終點（如死亡率、症狀改善）與次要終點（如生活品質、副作用）。例如：「心血管事件發生率」或「糖化血色素（HbA1c）下降幅度」。
\end{itemize}

PICO框架的應用可依臨床問題類型調整，例如診斷性問題可改為PICO的變體（Patient, Index test, Comparator test, Outcome），預後問題則可能省略I與C，聚焦於P與O。

透過PICO框架，模糊的臨床疑問可轉化為具體、可搜尋的問題，例如：
\begin{quote}
    非結構化問題：「運動對糖尿病患者有幫助嗎？」\\
    PICO結構化問題：「在患有2型糖尿病的成人患者（P）中，有氧運動（I）相較於標準藥物治療（C），是否能有效降低心血管事件風險（O）？」
\end{quote}

此框架不僅提升問題明確性，也直接對應搜尋策略的關鍵詞選擇。

\section{臨床問題的分類}

臨床問題依其性質可分為幾大類別，每類別對應不同的研究設計與證據類型。常見分類包括背景性問題（background questions）與前景性問題（foreground questions），後者更適用於實證醫學實踐。

主要臨床問題類型如下：
\begin{itemize}
    \item \textbf{治療/介入（Therapy/Intervention）}：探討某介入措施的有效性與安全性。典型問題：「哪種治療對某疾病最有效？」最適研究設計：隨機對照試驗（RCT）與系統回顧。
    \item \textbf{診斷（Diagnosis）}：評估診斷測試的準確性與可靠性。典型問題：「哪種測試最能準確診斷某疾病？」最適研究設計：橫斷面研究或診斷準確性研究。
    \item \textbf{預後（Prognosis）}：預測疾病的自然病程或長期結果。典型問題：「某疾病患者的長期存活率如何？」最適研究設計：縱向隊列研究（cohort study）。
    \item \textbf{病因/危害（Etiology/Harm）}：探討暴露因素與不良結果的因果關係。典型問題：「某暴露是否增加疾病風險？」最適研究設計：隊列研究或病例對照研究。
    \item \textbf{預防（Prevention）}：評估預防措施的有效性。可視為治療類型的延伸。
    \item \textbf{其他類型}：包括成本效益、患者體驗（質性研究）與流行頻率（prevalence）。
\end{itemize}

正確分類臨床問題有助於選擇適當的PICO變體與證據來源，例如治療問題強調RCT，而診斷問題則需敏感度與特異度等指標。

\section{問題轉化為可搜尋的查詢}

將結構化臨床問題轉化為高效的文獻搜尋查詢是實證醫學實踐的關鍵步驟。良好的查詢需平衡敏感度（召回率）與特異度（精準度）。

轉化步驟如下：
\begin{enumerate}
    \item 從PICO元件提取關鍵詞與同義詞。例如，P：「2型糖尿病」可擴展為「type 2 diabetes mellitus」或「NIDDM」。
    \item 使用醫學主題詞（MeSH）與自由文字詞結合。在PubMed等資料庫中，MeSH詞可提升檢索效率。
    \item 運用布林運算子（AND、OR、NOT）組合關鍵詞：
    \begin{itemize}
        \item OR：用於同義詞或相關概念（擴大範圍）。
        \item AND：用於不同PICO元件（縮小範圍）。
        \item NOT：排除不相關概念。
    \end{itemize}
    \item 加入限制條件（如出版年限、語言、研究類型）與過濾器（filter），例如PubMed的「Systematic Reviews」或「Randomized Controlled Trial」。
\end{enumerate}

範例查詢（PubMed格式）：
\begin{verbatim}
("Diabetes Mellitus, Type 2"[MeSH] OR "type 2 diabetes") 
AND ("Exercise"[MeSH] OR "aerobic exercise" OR "physical activity") 
AND ("Cardiovascular Diseases"[MeSH] OR "cardiovascular events" OR "myocardial infarction") 
AND ("Randomized Controlled Trial"[ptyp])
\end{verbatim}

此查詢對應前述PICO範例，結合MeSH詞、自由文字與研究類型過濾，可有效檢索相關高品質證據。

實務上，建議從敏感度高的廣泛搜尋開始，逐步精煉以獲得可管理的結果數量。熟悉資料庫介面與進階功能將顯著提升搜尋效率。

% 可在此新增表格（如PICO範例表）、圖示或習題
% \bibliographystyle{plain}
% \bibliography{references} % 若需參考文獻，請建立references.bib檔案

\end{document}

% 在此新增其他章節
% \include{chapter2}
% \include{chapter3}
% ...

\end{document}

實證醫學教材目錄設計建議
作為一個關於實證醫學（Evidence-Based Medicine, EBM）的教材，目錄的設計應以邏輯進階為原則，從基礎概念出發，逐步深入到方法論、應用實務，以及進階挑戰。這有助於讀者逐步建立知識框架，確保教材適合醫學學生、臨床醫師或研究人員使用。我的設計考量了EBM的核心元素，包括證據的產生、評估、整合與應用，同時融入倫理、統計和實踐指南，以維持全面性和實用性。目錄可分為主要章節，每章包含子節，以利於教學模組化。
以下是建議的教材目錄結構，共分為五個主要部分，涵蓋約10-15章，視教材深度而定。每章可附以學習目標、案例研究和討論題，以強化互動性。
第一部分：實證醫學基礎
此部分介紹EBM的核心概念和歷史背景，為後續內容奠定基礎。

實證醫學概述
定義與原則
EBM的歷史發展
EBM在醫療決策中的角色

醫療問題的制定
PICO框架（Patient, Intervention, Comparison, Outcome）
臨床問題的分類（診斷、治療、預後等）
問題轉化為可搜尋的查詢


第二部分：證據的產生與研究設計
此部分聚焦於科學研究的基礎方法，幫助讀者理解證據的來源。

研究設計類型
觀察性研究（cohort、case-control）
實驗性研究（隨機對照試驗，RCT）
質性研究與混合方法

統計學基礎
描述性與推斷性統計
置信區間、p值與統計顯著性
偏差與混淆因素

倫理與研究品質
研究倫理原則（Helsinki宣言）
證據等級體系（GRADE系統）
研究報告標準（CONSORT、STROBE）


第三部分：證據的搜尋與評估
此部分強調如何獲取和批判性分析證據，培養讀者的批判思維。

文獻搜尋策略
資料庫使用（PubMed、Cochrane Library）
搜尋術語與布林運算子
灰色文獻與開放存取資源

批判性評估
評估研究的效度、可靠性和適用性
工具如CASP checklist
偏差來源（選擇偏差、出版偏差）

系統回顧與meta-analysis
系統回顧的步驟（PRISMA指南）
Meta-analysis的統計方法
異質性與敏感度分析


第四部分：證據的整合與應用
此部分橋接理論與實務，指導如何將證據應用於臨床情境。

臨床指南與決策
指南開發過程（AGREE工具）
證據整合到患者照護
共享決策模型

EBM在專科領域的應用
診斷與篩檢
治療與預防
公共衛生與政策

實施EBM的障礙與策略
常見挑戰（時間限制、資源不足）
知識轉譯與實施科學
案例分析與模擬


第五部分：進階主題與未來展望
此部分探討EBM的擴展與前沿議題，適合進階讀者。

大數據與人工智慧在EBM中的角色
電子健康記錄與實時證據
機器學習在預測模型中的應用
倫理考量

EBM的全球視角
發展中國家的挑戰
文化與資源差異
國際合作與知識分享

未來方向與持續學習
新興趨勢（如精準醫學）
終身學習資源
評估EBM實踐的指標